%A palavra Sumário deve ser centralizada e com o mesmo tipo de fonte utilizada para as seções primárias. Traz as partes de que é composto o trabalho na ordem em que aparecem, seguidas da numeração das páginas. A formatação dos itens deverá ser na mesma tipologia e tamanho de fonte do texto. As partes devem ser alinhadas pelo indicativo mais extenso. Os elementos pré-textuais (capa, folha de rosto, resumo, listas, etc.) NÃO devem constar no sumário. Deve ser elaborado segundo a ABNT NBR 6027 (ASSOCIAÇÃO..., 2012b). Exemplo:}\\

\pdfbookmark[0]{\contentsname}{toc}
\tableofcontents*
\cleardoublepage